% !PRINTED 520
\weberpage{520}
Sind die gegebenen Functionale wirkliche Brüche, d. h. reduciren sie sich nicht alle auf ganze Functionale, so muss der gemeinsame Nenner wenigstens einen Primfactor enthalten, der nicht in allen Zählern vorkommt.

Alles das ist in vollkommener Uebereinstimmung mit den Regeln der elementaren Arithmetik, und auch die Beweismethoden, die wir hier angewandt haben, sind wesentlich dieselben, die dort gebraucht werden. Der Kernpunkt der Deduction ist einerseits die weitgehende Verallgemeinerung des Begriffes der Einheit, andererseits die darauf gegründete Definition des grössten gemeinschaftlichen Theilers im \S.\,139.

\begin{center}
\S.\,142.\\
\gesperrt{Ganze Functionen im Körper $\overline{\Omega}$.}
\end{center}

Die Hülfsmittel, über die wir jetzt verfügen, reichen aus, um die Schlussweise, deren wir uns im \S.\,2 des ersten Bandes zum Beweis des \gesperrt{Gauss}'schen Satzes bedient haben, auf die Functionale anzuwenden.

In der That ist ja der Satz, auf den sich jener Beweis wesentlich stützt, dass ein Product zweier ganzer Zahlen nur dann durch eine Primzahl theilbar ist, wenn diese Primzahl in einem der Factoren aufgeht, auch für die jetzt eingeführten Primfunctionale als gültig erwiesen.

Wir können demnach, indem wir dem erwähnten Beweise Schritt für Schritt folgen, den Satz als erwiesen ansehen:

1. \gesperrt{Wenn zwei ganze Functionen einer Variablen $t$ der Grade $h$ und $k$:}
\[
\varphi = \alpha_0 t^h + \alpha_1 t^{h-1} + \cdots + \alpha_h
\]
\[
\psi = \beta_0 t^k + \beta_1 t^{k-1} + \cdots + \beta_k,
\]
\gesperrt{deren Coëfficienten ganze Functionale sind, die die Variable $t$ nicht enthalten, ein Product}
\[
\chi = \gamma_0 t^{h+k} + \gamma_1 t^{h+k-1} + \cdots + \gamma_{h+k}
\]
\gesperrt{haben, in dem die Coëfficienten $\gamma_0, \gamma_1, \ldots, \gamma_{h+k}$ einen gemeinschaftlichen Primtheiler $\pi$ haben, so muss $\pi$ entweder in allen Coëfficienten von $\varphi$ oder in allen Coëfficienten von $\psi$ aufgehen.}
